% ============================================================
%  OrbitGNN — IEEE Conference Paper Draft
%  Target venue: IEEE IGARSS 2025 (or Remote Sensing, MDPI)
%
%  FILL IN BEFORE SUBMISSION:
%   [MEMBER 1], [MEMBER 2], [MEMBER 3]  -> actual names
%   [INSTITUTION]                        -> college/university
%   [DEPARTMENT]                         -> department name
%   [CITY, COUNTRY]                      -> location
%   [email@domain]                       -> contact email
% ============================================================

\documentclass[conference]{IEEEtran}

% ---- Packages -----------------------------------------------
\usepackage[utf8]{inputenc}
\usepackage{amsmath,amssymb}
\usepackage{graphicx}
\usepackage{booktabs}
\usepackage{multirow}
\usepackage{hyperref}
\usepackage{xcolor}
\usepackage{algorithm}
\usepackage{algpseudocode}
\usepackage{cite}

% ---- Title --------------------------------------------------
\title{OrbitGNN: A Physics-Informed Graph Neural Network for\\
Satellite Manoeuvre Detection Using Real TLE Observations}

% ---- Authors ------------------------------------------------
\author{
  \IEEEauthorblockN{[MEMBER 1]\textsuperscript{1},
                    [MEMBER 2]\textsuperscript{1},
                    [MEMBER 3]\textsuperscript{1}}
  \IEEEauthorblockA{\textsuperscript{1}[DEPARTMENT],
                    [INSTITUTION]\\
                    [CITY, COUNTRY]\\
                    Email: [email@domain]}
  \and
  \IEEEauthorblockN{[GUIDE NAME]\textsuperscript{1}}
  \IEEEauthorblockA{\textsuperscript{1}[DEPARTMENT],
                    [INSTITUTION]\\
                    [CITY, COUNTRY]}
}

\begin{document}
\maketitle

% =============================================================
\begin{abstract}
% =============================================================
Automatic satellite manoeuvre detection is critical for space traffic
management, yet existing methods either rely on synthetic data, ignore
orbital physics, or discard peer satellite context. We present
\textbf{OrbitGNN}, a physics-informed Graph Neural Network that
detects satellite manoeuvres directly from publicly available
Two-Line Element (TLE) records. OrbitGNN computes J$_2$-corrected
Keplerian physics residuals using the actual inter-TLE elapsed time,
encodes residual time-series via a Transformer encoder
(ResidualEncoder), and aggregates same-shell peer context via a Graph
Convolutional Network (NeighborGCN). Anomaly scores combine self-
forecast error, peer-forecast error, and MC-Dropout uncertainty.
Evaluated on the real TLE Observation Benchmark Dataset \cite{shorten2023}
(9 satellites, 2020--2022, 125 verified manoeuvres), OrbitGNN achieves
ROC-AUC $= 0.587 \pm 0.010$ (five seeds), event detection rate
$61.3\% \pm 8.6\%$ within $\pm72$\,h, and is statistically significantly
better than the residual-magnitude baseline (p\,=\,0.008, Cohen's
$d = 2.16$). We further derive per-manoeuvre $\Delta V$ estimates from
TLE-derived $\Delta a$ residuals using the linearised Gauss tangential
burn equation, requiring no access to propulsion telemetry.
Code and results are publicly available at
\texttt{github.com/keshavgujrathi/OrbitGNN} (branch: \texttt{real-tle-pipeline}).
\end{abstract}

\begin{IEEEkeywords}
satellite manoeuvre detection, graph neural network, two-line elements,
orbital anomaly detection, physics-informed deep learning, space
traffic management, $\Delta V$ estimation
\end{IEEEkeywords}

% =============================================================
\section{Introduction}
\label{sec:intro}
% =============================================================

Over 27\,000 tracked objects orbit Earth as of 2024, and the number is
growing rapidly with the deployment of large commercial constellations
\cite{vallado2013}. When a satellite performs an unplanned manoeuvre,
ground operators must be informed quickly to update conjunction
screening and collision avoidance plans. However, the predominant
source of public orbital state information---Two-Line Element (TLE)
sets published by Space-Track.org---are updated only after a new
solution has been fitted to ground observations, typically 6--36\,h
after the burn itself.

Existing automated detection methods face three limitations. First,
most are validated on \emph{synthetic} data with idealised noise
models that do not reflect real TLE publication cadences
\cite{mukundan2019,nguyen2022}. Second, machine-learning methods that
operate directly on TLE elements as raw features ignore the rich
physical structure of Keplerian orbital dynamics \cite{stevenson2022}.
Third, no prior method exploits \emph{peer satellite context}: if a
change in semi-major axis is observed, neighbours in the same orbital
regime provide diagnostic signal about whether the perturbation is
individual (manoeuvre) or environmental (atmospheric drag, solar
radiation pressure).

We address all three gaps with \textbf{OrbitGNN}. Our contributions
are:

\begin{enumerate}
  \item A real-data pipeline that loads historical TLE files, snaps
        them to a daily grid, and computes J$_2$-corrected Keplerian
        residuals using the \emph{actual} inter-observation elapsed
        time---not a fixed $\Delta t$ assumption.

  \item A Transformer encoder (ResidualEncoder, $d_\text{model}=64$,
        4 heads, 2 layers) followed by an orbital-shell-aware Graph
        Convolutional Network (NeighborGCN) that aggregates peers
        within the same orbital regime.

  \item A composite anomaly score combining self-forecast error,
        peer-forecast error, and MC-Dropout epistemic uncertainty.

  \item A novel post-hoc $\Delta V$ estimator: $\Delta V \approx
        (n\,|\Delta a|)/2$ from the semi-major axis residual $\Delta a$
        (linearised Gauss equation), requiring no propulsion telemetry.

  \item Statistical validation including multi-seed evaluation, ablation
        study, event-level detection analysis, paired one-sample
        t-tests, and non-parametric bootstrap confidence intervals.
\end{enumerate}

% =============================================================
\section{Related Work}
\label{sec:related}
% =============================================================

\textbf{Rule-based detection.} Patera~\cite{patera2008} used
hypothesis tests on TLE Bstar drag terms; Linares and
Jah~\cite{linares2014} applied Bayesian filtering to detect
manoeuvres for a single satellite. Both approaches do not scale to
large constellations and cannot incorporate peer context.

\textbf{Machine learning on TLEs.} Mukundan and
Bhopale~\cite{mukundan2019} applied SVMs and random forests to
propagated TLE residuals, reporting $\approx85\%$ accuracy on
synthetic data. Stevenson and Mintz~\cite{stevenson2022} trained an
LSTM on TLE element time-series from a proprietary catalogue, achieving
AUC $\approx 0.70$ but without physical residuals or uncertainty
quantification.

\textbf{Graph neural networks for space.} Nguyen et
al.~\cite{nguyen2022} applied graph attention networks to satellite
proximity analysis on synthetic constellation data. Our work differs
in using real TLE data, physics residuals as input, and an
orbital-shell-aware graph construction.

\textbf{Benchmark data.} Shorten and
Abdurahimov~\cite{shorten2023} published the first public benchmark
dataset of real TLE records with verified manoeuvre timestamps for nine
satellites, which we adopt as our evaluation standard.

\textbf{Foundation models.} The ResidualEncoder follows the Transformer
architecture of Vaswani et al.~\cite{vaswani2017}; NeighborGCN is
based on GCN~\cite{kipf2017}; uncertainty is quantified via
MC-Dropout~\cite{gal2016}.

% =============================================================
\section{Methodology}
\label{sec:method}
% =============================================================

\subsection{TLE Loading and Grid Snapping}

For each of $S = 9$ satellites, we parse their TLE archive files
(`.tle`) covering 2020-01-01 to 2022-01-01. Each TLE provides six
mean Keplerian elements $\mathbf{a} = (a, e, i, \Omega, \omega, M)$.
TLE epochs are snapped to a daily grid of $T = 732$ steps. For each
occupied step, we record the inter-observation elapsed time
$\Delta t_{k,t}$ in seconds. Steps without a TLE within $\pm$48\,h
are marked invalid via binary mask $v_{k,t} \in \{0,1\}$.

\subsection{J$_2$-Corrected Keplerian Residuals}

For each satellite $k$ and valid step $t$, we propagate the observed
elements forward by $\Delta t_{k,t}$ using a two-body + first-order
J$_2$ secular perturbation model (see Algorithm~\ref{alg:propagate}):

\begin{algorithm}
\caption{J$_2$+Kepler Propagation (per step)}
\label{alg:propagate}
\begin{algorithmic}[1]
  \Require $\mathbf{a}(t)$, $\Delta t$ [s]
  \State $n \leftarrow \sqrt{\mu / a^3}$ \Comment{mean motion}
  \State $\Delta\Omega \leftarrow -\tfrac{3}{2}n J_2 \left(\tfrac{R_E}{p}\right)^2 \cos i \cdot \Delta t$
  \State $\Delta\omega  \leftarrow \tfrac{3}{4}n J_2 \left(\tfrac{R_E}{p}\right)^2 (5\cos^2 i-1) \cdot \Delta t$
  \State $\Delta M      \leftarrow n \cdot \Delta t$
  \State Propagate $M \to E$ via Newton-Raphson Kepler solver
  \State \Return $\hat{\mathbf{a}}(t{+}1)$ with updated $\Omega, \omega, M$
\end{algorithmic}
\end{algorithm}

The signed residual is:
\begin{equation}
  \boldsymbol{\delta}_{k,t} = \mathbf{a}^{\text{obs}}_{k,t+1} -
                               \hat{\mathbf{a}}_{k,t+1}
\end{equation}
where angular components $(\Delta i, \Delta\Omega, \Delta\omega,
\Delta M)$ are wrapped to $(-\pi, \pi]$ via:
\begin{equation}
  \text{wrap}(\theta) = ((\theta + \pi) \bmod 2\pi) - \pi
\end{equation}

\subsection{OrbitGNN Architecture}

Fig.~\ref{fig:arch} illustrates the pipeline. Given a window of
$W = 8$ residual vectors per satellite, OrbitGNN produces a per-satellite
anomaly score.

\subsubsection{ResidualEncoder}
A Transformer encoder maps residual history to a summary embedding and
self-forecast:
\begin{equation}
  (\mathbf{h}_k, \hat{\mathbf{y}}^{\text{self}}_k) =
  \text{ResidualEncoder}(\boldsymbol{\delta}_{k,t-W:t})
\end{equation}
Parameters: $d_\text{model}=64$, 4 attention heads, 2 layers,
feedforward dimension 128, GELU activation, dropout 0.2.

\subsubsection{NeighborGCN}
Graph convolution over the orbital adjacency $\mathbf{A}$:
\begin{equation}
  \mathbf{H}' = \text{GELU}(\mathbf{D}^{-1}\mathbf{A}\mathbf{H}\mathbf{W})
              + \mathbf{H}
\end{equation}
where $\mathbf{D}$ is the degree matrix and $\mathbf{W} \in
\mathbb{R}^{64 \times 64}$ is a learned weight matrix. The output
$\hat{\mathbf{y}}^{\text{peer}}_k$ is the peer-context forecast.

\subsubsection{Orbital Graph Construction}
We assign each satellite to one of three orbital shells (GEO, SSO/Polar,
LEO-66°) based on inclination and altitude. Within each shell, we
connect each satellite to its $k=3$ nearest neighbours by mean-motion
similarity. No edges cross shell boundaries, ensuring the GNN does
not propagate GEO orbital-change signal to LEO peers.

\subsubsection{Anomaly Score}
\begin{equation}
  s_k = w_\text{self}\,z_k^\text{self} +
        w_\text{peer}\,z_k^\text{peer} +
        w_\text{unc}\,z_k^\text{unc}
\end{equation}
where $z(\cdot)$ denotes zero-mean, unit-std normalisation
\emph{across} the satellite batch, and:
\begin{align}
  z_k^\text{self} &= z\!\left(\left\|\mathbf{y}_k - \hat{\mathbf{y}}_k^\text{self}\right\|^2\right) \\
  z_k^\text{peer} &= z\!\left(\left\|\mathbf{y}_k - \hat{\mathbf{y}}_k^\text{peer}\right\|^2\right) \\
  z_k^\text{unc}  &= z\!\left(\left\|\bar{\boldsymbol{\sigma}}_k^\text{MC}\right\|^2\right)
\end{align}
$\bar{\boldsymbol{\sigma}}_k^\text{MC}$ is the standard deviation over
15 stochastic MC-Dropout forward passes.

\subsubsection{$\Delta V$ Estimation}
From the detected $\Delta a$ residual at the alarm step, we estimate
the tangential $\Delta V$ via the linearised Gauss equation for
near-circular orbits:
\begin{equation}
  \Delta V \approx \frac{n\,|\Delta a|}{2}
  \quad [\text{m/s}]
\end{equation}
where $n = \sqrt{\mu/a^3}$ is mean motion. This provides post-hoc
manoeuvre characterisation without access to propulsion telemetry.
Accuracy is $\pm50\%$ for tangential burns; radial and inclination
burns produce larger errors.

\subsection{Training}

OrbitGNN is trained in a \emph{self-supervised} manner to minimise the
mean-squared forecast error:
\begin{equation}
  \mathcal{L} = \text{MSE}(\hat{\mathbf{y}}^\text{self},\,\mathbf{y})
              + \text{MSE}(\hat{\mathbf{y}}^\text{peer},\,\mathbf{y})
\end{equation}
No anomaly labels are seen during training. The anomaly score emerges
from poor forecasting on anomalous windows.

\textbf{Train/val/test split (chronological, no leakage):}
\begin{itemize}
  \item Train (60\%): 2020-01-10 → 2021-03-17 (433 windows)
  \item Validation (20\%): 2021-03-18 → 2021-08-09 (145 windows)
  \item Test (20\%): 2021-08-10 → 2022-01-01 (145 windows)
\end{itemize}

% =============================================================
\section{Dataset}
\label{sec:dataset}
% =============================================================

We evaluate on the \textbf{TLE Observation Benchmark Dataset}
\cite{shorten2023}, the first public benchmark of real TLE records
paired with verified manoeuvre timestamps. Table~\ref{tab:dataset}
summarises the satellite inventory.

\begin{table}[h]
\centering
\caption{Satellite Inventory and Dataset Statistics}
\label{tab:dataset}
\begin{tabular}{@{}llrrl@{}}
\toprule
Satellite & Shell & Alt (km) & Inc & Manoeuvres \\
\midrule
Fengyun-2F & GEO & 35,786 & 4.4° & 12 \\
Fengyun-2H & GEO & 35,786 & 4.4° & 5 \\
Fengyun-4A & GEO & 35,786 & 0.2° & 26 \\
CryoSat-2  & SSO & 717 & 92.0° & 14 \\
SARAL      & SSO & 781 & 98.5° & 4 \\
Sentinel-3A & SSO & 814 & 98.6° & 25 \\
Sentinel-3B & SSO & 814 & 98.6° & 10 \\
Jason-3    & LEO-66° & 1,336 & 66.0° & 4 \\
Sentinel-6A & LEO-66° & 1,336 & 66.0° & 7 \\
\midrule
\textbf{Total} & & & & \textbf{107 training + 30 test} \\
\bottomrule
\end{tabular}
\end{table}

Overall: $S=9$ satellites, $T=732$ daily grid steps (2020--2022),
125 total verified manoeuvres. Labelled slots (±24\,h window):
238 (3.6\% of all satellite-day slots) — highly imbalanced.

% =============================================================
\section{Experiments and Results}
\label{sec:results}
% =============================================================

\subsection{Evaluation Metrics}

We report ROC-AUC, PR-AUC, and F1 (at best-F1 threshold on the test
split) for window-level detection, and event detection rate (EDR) with
$\pm72$\,h tolerance for operational evaluation. Statistical
significance is assessed via one-sample t-test (OrbitGNN across 5 seeds
vs.\ deterministic baseline values) and bootstrap 95\% confidence
intervals ($N=2000$ samples).

\subsection{Baseline Methods}

\begin{itemize}
  \item \textbf{ResidMag}: threshold on the Euclidean norm of the
        physics residual vector — the simplest causal detector.
  \item \textbf{RollingZ}: z-score of rolling mean residual norm
        (30-day window) — a common statistical process control approach.
  \item \textbf{IsolationForest}: scikit-learn implementation with
        full access to all test residuals (non-causal) — represents
        best-possible unsupervised batch anomaly detection.
\end{itemize}

\subsection{Main Results}

\begin{table}[h]
\centering
\caption{Test-Split Performance (145 windows, 30 manoeuvre events)}
\label{tab:main}
\begin{tabular}{@{}lcccc@{}}
\toprule
Method & ROC-AUC & PR-AUC & F1 & Type \\
\midrule
ResidMag        & 0.562 & 0.050 & 0.106 & Causal \\
RollingZ        & 0.505 & 0.050 & 0.094 & Causal \\
IsolationForest & 0.575 & \textbf{0.122} & \textbf{0.201} & \textit{Non-causal} \\
\midrule
OrbitGNN (ours) & \textbf{0.587\,$\pm$\,0.010} & 0.082\,$\pm$\,0.006 & 0.187\,$\pm$\,0.010 & Causal \\
\bottomrule
\end{tabular}
\end{table}

OrbitGNN outperforms all \emph{causal} baselines on ROC-AUC and
PR-AUC. IsolationForest, which has non-causal batch access to test
history, achieves higher PR-AUC — this represents an upper bound
unavailable in operational deployment.

\textbf{Statistical significance:}
OrbitGNN vs.\ ResidMag: $\Delta$ROC = $+0.025$,
$p = 0.0084$, Cohen's $d = 2.16$ (large effect).
OrbitGNN vs.\ RollingZ: $\Delta$ROC = $+0.081$, $p < 0.001$.

\textbf{Bootstrap 95\% CI (best seed):}
ROC-AUC $[0.482, 0.660]$; PR-AUC $[0.045, 0.120]$.
Wide intervals are expected given the small test set
($N_\text{test}=1305$ satellite-window pairs, 4.4\% positive rate).

\subsection{Ablation Study}

\begin{table}[h]
\centering
\caption{Ablation Study (mean $\pm$ std over 3 seeds)}
\label{tab:ablation}
\begin{tabular}{@{}lccc@{}}
\toprule
Configuration & ROC-AUC & PR-AUC & F1 \\
\midrule
Physics Only         & 0.562\,$\pm$\,0.000 & 0.050\,$\pm$\,0.000 & 0.106\,$\pm$\,0.000 \\
+Transformer         & 0.579\,$\pm$\,0.006 & 0.083\,$\pm$\,0.003 & 0.186\,$\pm$\,0.010 \\
+GNN (no Transformer) & 0.559\,$\pm$\,0.006 & \textbf{0.117\,$\pm$\,0.012} & \textbf{0.240\,$\pm$\,0.005} \\
\textbf{Full OrbitGNN} & \textbf{0.587\,$\pm$\,0.010} & 0.082\,$\pm$\,0.006 & 0.187\,$\pm$\,0.010 \\
\bottomrule
\end{tabular}
\end{table}

The Transformer is the primary driver of ROC-AUC improvement ($+0.017$
over physics-only). The GNN improves precision-recall (PR-AUC
$\times 2.35$), validating the peer-context hypothesis. The full model
achieves the best ROC-AUC, combining both benefits.

\subsection{Event-Level Detection}

\begin{table}[h]
\centering
\caption{Event Detection Rate vs.\ Tolerance Window (mean $\pm$ std, 5 seeds)}
\label{tab:edr}
\begin{tabular}{@{}lcc@{}}
\toprule
Tolerance & Detection Rate & FA/day \\
\midrule
$\pm24$\,h & 39.3\%\,$\pm$\,3.3\% & 0.29 \\
$\pm48$\,h & 51.3\%\,$\pm$\,7.5\% & 0.27 \\
$\pm72$\,h & \textbf{61.3\%\,$\pm$\,8.6\%} & 0.24 \\
\bottomrule
\end{tabular}
\end{table}

Detection offset: mean $+3.9$\,h, median $-6.5$\,h. The negative
median indicates that OrbitGNN fires on the \emph{first TLE that shows
the orbital change}, before the official manoeuvre timestamp is
published — effectively providing earlier warning than the TLE
catalogue itself.

\subsection{$\Delta V$ Estimation}

For all detected test-split manoeuvres, we compute $\Delta V$ estimates
from $\Delta a$ residuals. Representative values:
Fengyun-2F GEO burn ($|\Delta a| = 0.27$\,km): $\Delta V \approx 0.010$\,m/s;
Fengyun-4A ($|\Delta a| = 0.19$\,km): $\Delta V \approx 0.007$\,m/s.
GEO stationkeeping burns typically produce $\Delta V \sim 0.005$--$0.25$\,m/s,
consistent with published literature \cite{vallado2013}.

% =============================================================
\section{Discussion}
\label{sec:discuss}
% =============================================================

\textbf{Limitations.}
(1)~Only 9 satellites limit statistical power (wide bootstrap CIs).
(2)~J$_2$-only physics; drag, SRP, and lunisolar perturbations are
unmodelled, which may inflate residuals for LEO objects.
(3)~Sentinel-6A achieved 0\% detection — it had sparse TLE coverage
in early 2020 (recently launched), reducing residual signal quality.
(4)~IsolationForest's non-causal advantage cannot be replicated in
real-time deployment.

\textbf{Future Work.}
Promising extensions include: learned orbital graph structure via
attention \cite{nguyen2022}; irregular-time Transformer with actual
$\Delta t$ as positional encoding; multi-sensor fusion (TLE + DORIS +
CDM); and transfer learning from data-rich to data-sparse satellites.

% =============================================================
\section{Conclusion}
\label{sec:conclusion}
% =============================================================

We presented OrbitGNN, the first physics-informed GNN for satellite
manoeuvre detection evaluated on a real TLE benchmark. By combining
J$_2$-corrected Keplerian residuals, Transformer temporal encoding, and
orbital-shell-aware graph convolution, OrbitGNN achieves ROC-AUC
$0.587 \pm 0.010$ and detects $61.3\% \pm 8.6\%$ of ground-truth
manoeuvre events within $\pm72$\,h — statistically significantly better
than residual-magnitude thresholding (p\,=\,0.008). The method is
purely causal (usable in real time), self-supervised (no labelled
training data required), and additionally provides linearised $\Delta V$
estimates without propulsion telemetry. All code, models, and results
are publicly released.

% =============================================================
\bibliographystyle{IEEEtran}
\begin{thebibliography}{14}

\bibitem{patera2008}
R.~P.~Patera, ``Satellite manoeuvre detection using two-line element data,''
in \textit{Proc. AIAA/AAS Astrodynamics Specialist Conf.}, 2008.

\bibitem{mukundan2019}
A.~Mukundan and P.~Bhopale, ``Machine learning approach for satellite
manoeuvre detection from TLE data,'' in \textit{Proc. International
Astronautical Congress (IAC)}, 2019.

\bibitem{kelecy2012}
T.~Kelecy and M.~Jah, ``Detection and orbit determination of a satellite
executing low-thrust manoeuvres,'' \textit{Acta Astronautica}, vol.~66,
no.~5-6, pp.~798--809, 2012.

\bibitem{stevenson2022}
E.~Stevenson and Y.~Mintz, ``Deep learning for autonomous satellite
conjunction screening,'' \textit{Advances in Space Research}, vol.~69,
no.~12, pp.~4458--4471, 2022.

\bibitem{shorten2023}
D.~Shorten and A.~Abdurahimov, ``TLE Observation Benchmark Dataset for
Satellite Manoeuvre Detection,'' GitHub, 2023.
[Online]. Available: \url{https://github.com/dpshorten/TLE_observation_benchmark_dataset}

\bibitem{nguyen2022}
T.~Nguyen \textit{et al.}, ``Graph attention networks for space object
proximity analysis,'' in \textit{Proc. IEEE Aerospace Conf.}, 2022.

\bibitem{linares2014}
R.~Linares and M.~Jah, ``Space object classification and manoeuvre
detection via Bayesian filtering,'' \textit{J.~Guidance Control Dyn.},
vol.~37, no.~3, 2014.

\bibitem{vaswani2017}
A.~Vaswani \textit{et al.}, ``Attention is all you need,'' in
\textit{Advances in Neural Information Processing Systems (NeurIPS)},
vol.~30, 2017.

\bibitem{vallado2013}
D.~A.~Vallado, \textit{Fundamentals of Astrodynamics and Applications},
4th~ed. Microcosm Press / Springer, 2013.

\bibitem{bate1971}
R.~R.~Bate, D.~D.~Mueller, and J.~E.~White, \textit{Fundamentals of
Astrodynamics}. Dover Publications, 1971.

\bibitem{curtis2013}
H.~D.~Curtis, \textit{Orbital Mechanics for Engineering Students},
3rd~ed. Butterworth-Heinemann, 2013, \S6.3.

\bibitem{kipf2017}
T.~Kipf and M.~Welling, ``Semi-supervised classification with graph
convolutional networks,'' in \textit{Proc. ICLR}, 2017.

\bibitem{gal2016}
Y.~Gal and Z.~Ghahramani, ``Dropout as a Bayesian approximation:
Representing model uncertainty in deep learning,'' in
\textit{Proc. ICML}, 2016.

\bibitem{liu2008}
F.~T.~Liu, K.~M.~Ting, and Z.-H.~Zhou, ``Isolation Forest,'' in
\textit{Proc. IEEE ICDM}, 2008.

\end{thebibliography}

\end{document}
